\chapter{Prolegomena to Finding Mind}
\label{chap:prolegomena}

If this book is going to be written, it has to come on its own terms. It's living in a mind that is as much musician and poet as mathematician and computer scientist. So, i'm going to let go of the traditional forms one might normally employ to grapple with content like this. i'm going to let this narrative be as personal as it needs to be or as impersonal as it needs to be to create the conditions whereby the story can be told.

What story? Well, part of the story is about a relationship between origin of life and origin of mind. In a nutshell, the idea is that origin of mind is a recapitulation of the origin of life, one octave up. (There's the musician talking.) In order to tell that story, however i need to retell the story of what a computer is. That retelling will need to address a peculiar blind spot that plagues us all. More on that later.

What's the connection between life and mind and computers? Information. You see everyone is somehow convinced that computational and/or information theoretic approaches have an extraordinary explanatory power. And i mean everyone from biologists like David Sinclair who is involved in clinical trials testing an especially effective theory of aging as an information theoretic phenomenon, to physicists working at the information theoretic interface of black holes and quantum mechanics to AI researchers who are building information theoretic models of human cognition. It only stands to reason that we ought to explore an information theoretic account of origin of life and origin of mind.

Indeed researchers like Lee Cronin and Sarah Walker are already exploring information theoretic approaches to origin of life. Their work on assembly theory offers keen insight into certain aspects of causation and memory that i believe play an important role in the story.

The thing about research into deep questions like aging or quantum gravity or artificial intelligence or origin of life is that it requires so much focus on the topic at hand that it is very difficult to keep up with developments in other fields. In particular, the story of what a computer is has come quite a long ways since Alan Turing's initial attempts to characterize computation. Thus, while these research programmes are using approaches informed by computation they are, like many people using actual computers, often using systems that are several versions behind.

So, part of the effort here is to get the readers of this book on the latest version of the story about what constitutes a computer and computation more generally. After the system update, so to speak, we can look at a bigger picture, one that addresses experimentally based learning, how that relates to logic, and how logic embeds into a scientifically driven exploration of one's environment and one's self.

Indeed, part of the story unfolding here is that the hard problem of consciousness lies not so much in the question ``What is it like to be a bat,'' as Thomas Nagel put it, but rather in the question, ``What is it like to be a computer program?'' 

The methodology i'm going to employ --- as if you hadn't already guessed --- is to write a collection of Substack articles, like this one, that i can assemble into the book. The title of the book is Finding Mind. The structure of the book is the structure of a magic trick: the \textbf{Pledge (showing the ordinary), the Turn (making it do the extraordinary), and the Prestige (bringing it back or revealing the impossible). }The first group of articles will constitute the Pledge. The next the Turn. The last the Prestige.

i have cancelled my streaming subscriptions for a year to make space to allow this book to happen. i'm hoping readers will come along to help me stay accountable to this aim.
